% ============================================================
%  Chapter 1 -- Exercises 1.1-1.20
% ============================================================
\rsection{Exercises}

\begin{roadmap}[How to use these]
Rudin's exercises are not drill. Several of them --- 1.6 and 1.7 in particular
--- prove results the book will later use, and skipping them leaves a hole.
The annotation here gives the \emph{approach}: which theorem to reach for, what
shape the argument takes, where the difficulty actually is. It does not give
solutions. Read the note only after you are genuinely stuck, and never before
attempting the problem.

\medskip
\raggedright
\textbf{The ones that matter most.} \emph{1.6 and 1.7} construct $b^x$ for real
$x$ and its inverse, the logarithm --- Rudin defers this to the exercises and
then uses exponentials freely from Chapter 3 onward. \emph{1.5} is the $\inf/\sup$
duality you will use constantly. \emph{1.8} explains why $\C$ is never given an
order. \emph{1.20} tests whether you understood why a cut may have no largest
member.

\medskip
\textbf{The routine ones.} 1.1--1.4 and 1.11--1.14 are short consequences of
1.14--1.18 and 1.33; do them quickly. 1.16--1.19 are geometry in $\R^k$ and
exercise the Schwarz inequality.
\end{roadmap}

\noindent
\emph{Unless the contrary is explicitly stated, all numbers that are mentioned
in these exercises are understood to be real.}

\begin{enumerate}[label=\textbf{1.\arabic*.}, leftmargin=3.0em, itemsep=9pt,
                  topsep=8pt, labelsep=0.5em]

\item If $r$ is rational ($r \ne 0$) and $x$ is irrational, prove that $r+x$ and
$rx$ are irrational.\kn{Contradiction, and the point is that $\Q$ is closed
under the operations. If $r+x$ were rational, then $x = (r+x)-r$ would be a
difference of two rationals, hence rational. Same for $rx$, dividing by $r$ ---
which is where $r \ne 0$ is needed.}

\item Prove that there is no rational number whose square is $12$.\kn{The parity
argument of 1.1 does not transfer directly, since $12$ is not prime. Work with
the factor $3$ instead of $2$: if $m^2 = 12n^2$ then $3 \mid m^2$, hence
$3 \mid m$ (this needs $3$ prime), and the descent proceeds as before.}

\item Prove Proposition 1.15.

\item Let $E$ be a nonempty subset of an ordered set; suppose $\alpha$ is a lower
bound of $E$ and $\beta$ is an upper bound of $E$. Prove that
$\alpha \le \beta$.\kw{The hypothesis ``nonempty'' is the entire exercise, and
it is easy to write a proof that never uses it --- such a proof is wrong. Pick
any $x \in E$; then $\alpha \le x \le \beta$. If $E$ were empty, every element
would be both a lower and an upper bound and the conclusion would fail.}

\item Let $A$ be a nonempty set of real numbers which is bounded below. Let $-A$
be the set of all numbers $-x$, where $x \in A$. Prove that\kn{Reflection swaps the two notions: $\beta$ is a lower bound of $A$ if and
only if $-\beta$ is an upper bound of $-A$, and the inequality reverses on
negation. Verify both clauses of Definition 1.8 for $-\sup(-A)$. This is a
second proof of 1.11 for the special case $S = \R$, and the identity is used
without comment throughout Chapters 3 and 6. It is Tao's Exercise 5.5.1.}
\[ \inf A = -\sup(-A). \]

\item Fix $b > 1$.\kn{This is a construction, not a computation, and it is the
exercise Rudin relies on most. The architecture: (a) shows the rational exponent
is well defined --- raise both sides to the power $nq$ and use uniqueness in
1.21; (b) is then algebra; (c) is the crucial step, showing the new definition
by supremum agrees with the old one on rationals, so that $b^x := \sup B(x)$
extends rather than contradicts; (d) is the payoff. Part (c) is where $b>1$
matters --- it makes $t \mapsto b^t$ increasing, so $B(r)$ has $b^r$ as its
largest element. If you get stuck, this construction is carried out in the main
text of Tao's \emph{Analysis I}, \S5.6 and \S6.7 --- what Rudin relegates to an
exercise, Tao develops as numbered propositions.}
  \begin{enumerate}[label=(\alph*), leftmargin=2.2em, itemsep=3pt, topsep=4pt]
    \item If $m,n,p,q$ are integers, $n>0$, $q>0$, and $r = m/n = p/q$, prove
          that
          \[ (b^m)^{1/n} = (b^p)^{1/q}. \]
          Hence it makes sense to define $b^r = (b^m)^{1/n}$.
    \item Prove that $b^{r+s} = b^rb^s$ if $r$ and $s$ are rational.
    \item If $x$ is real, define $B(x)$ to be the set of all numbers $b^t$,
          where $t$ is rational and $t \le x$. Prove that
          \[ b^r = \sup B(r) \]
          when $r$ is rational. Hence it makes sense to define
          \[ b^x = \sup B(x) \]
          for every real $x$.
    \item Prove that $b^{x+y} = b^xb^y$ for all real $x$ and $y$.
  \end{enumerate}

\item Fix $b>1$, $y>0$, and prove that there is a unique real $x$ such that
$b^x = y$, by completing the following outline. (This $x$ is called the
\emph{logarithm} of $y$ to the base $b$.)\kn{Compare the whole outline with the proof of 1.21 and the resemblance is
exact: form the set of things that undershoot, take its supremum, and rule out
both strict inequalities. Parts (a)--(e) are the continuity estimates that
Rudin computed by hand there --- (d) and (e) are the analogues of choosing
$h$ and $k$. Part (a) is the binomial theorem, or induction. Once you see that
1.21 and 1.7 are the same argument twice, the outline reads itself.}
  \begin{enumerate}[label=(\alph*), leftmargin=2.2em, itemsep=2pt, topsep=4pt]
    \item For any positive integer $n$, $b^n - 1 \ge n(b-1)$.
    \item Hence $b - 1 \ge n(b^{1/n}-1)$.
    \item If $t>1$ and $n > (b-1)/(t-1)$, then $b^{1/n} < t$.
    \item If $w$ is such that $b^w < y$, then $b^{w+(1/n)} < y$ for sufficiently
          large $n$; to see this, apply part (c) with $t = y\cdot b^{-w}$.
    \item If $b^w > y$, then $b^{w-(1/n)} > y$ for sufficiently large $n$.
    \item Let $A$ be the set of all $w$ such that $b^w < y$, and show that
          $x = \sup A$ satisfies $b^x = y$.
    \item Prove that this $x$ is unique.
  \end{enumerate}

\item Prove that no order can be defined in the complex field that turns it into
an ordered field. \emph{Hint:} $-1$ is a square.\kn{Two lines from 1.18(d).
That proposition gives $x^2>0$ for $x \ne 0$ in \emph{any} ordered field, so
$i^2 = -1 > 0$; but $1 > 0$ also, and then (i) of 1.17 forces
$0 = 1 + (-1) > 0$. The point is that no cleverness in choosing the order can
help --- the obstruction is 1.18(d), which holds for every order whatever.}

\item Suppose $z = a+bi$, $w = c+di$. Define $z < w$ if $a<c$, and also if $a=c$
but $b<d$. Prove that this turns the set of complex numbers into an ordered set.
(This type of order relation is called a \emph{dictionary order}, or
\emph{lexicographic order}, for obvious reasons.) Does this ordered set have the
least-upper-bound property?\kd{Read alongside 1.8. An order on $\C$ certainly
exists; what 1.8 rules out is an order \emph{compatible with the arithmetic}.
Definition 1.6 asks only for trichotomy and transitivity, and lexicographic
order supplies both. For the second question, test the set
$\{z : \operatorname{Re}(z)<0\}$ and ask what its least upper bound would have
to be.}

\item Suppose $z = a+bi$, $w = u+iv$, and
\[ a = \left(\frac{|w|+u}{2}\right)^{1/2}, \qquad
   b = \left(\frac{|w|-u}{2}\right)^{1/2}. \]
Prove that $z^2 = w$ if $v \ge 0$ and that $(\bar z)^2 = w$ if $v \le 0$.
Conclude that every complex number (with one exception!) has two complex square
roots.\kn{Direct computation: expand $z^2 = a^2-b^2+2abi$ and check the two
parts separately. The real part is immediate; for the imaginary part you will
need $|w|^2 = u^2+v^2$ and a square root taken with the right sign, which is
exactly where the hypothesis on $v$ enters. Both $a$ and $b$ are well defined
because $|w| \ge |u|$ by 1.33(d). The exception is $w = 0$.}

\item If $z$ is a complex number, prove that there exists an $r \ge 0$ and a
complex number $w$ with $|w| = 1$ such that $z = rw$. Are $w$ and $r$ always
uniquely determined by $z$?\kn{Take $r = |z|$ and $w = z/|z|$, then handle
$z = 0$ separately --- and that separate case is the answer to the second
question.}

\item If $z_1,\dots,z_n$ are complex, prove that\kn{Induction on $n$, with 1.33(e) as the base and the inductive step. Worth
doing once carefully: this is the form of the triangle inequality actually used
in Chapter 3.}
\[ |z_1+z_2+\cdots+z_n| \le |z_1|+|z_2|+\cdots+|z_n|. \]

\item If $x,y$ are complex, prove that\kn{The reverse triangle inequality. Write $x = (x-y)+y$ and apply 1.33(e) to
get $|x| - |y| \le |x-y|$; then swap $x$ and $y$ to get the other direction, and
combine. Used constantly later to show that the modulus is continuous.}
\[ \bigl||x|-|y|\bigr| \le |x-y|. \]

\item If $z$ is a complex number such that $|z| = 1$, that is, such that
$z\bar z = 1$, compute\kn{Expand each term as $(1\pm z)(\overline{1 \pm z})$ using 1.31, and watch the
cross terms cancel. The answer is a constant --- this is the parallelogram law
of 1.17 in disguise.}
\[ |1+z|^2 + |1-z|^2. \]

\item Under what conditions does equality hold in the Schwarz
inequality?

\item Suppose $k \ge 3$, $\mathbf{x},\mathbf{y} \in \R^k$,
$|\mathbf{x}-\mathbf{y}| = d > 0$, and $r > 0$. Prove:
  \begin{enumerate}[label=(\alph*), leftmargin=2.2em, itemsep=2pt, topsep=4pt]
    \item If $2r > d$, there are infinitely many $\mathbf{z} \in \R^k$ such that
          \[ |\mathbf{z}-\mathbf{x}| = |\mathbf{z}-\mathbf{y}| = r. \]
    \item If $2r = d$, there is exactly one such $\mathbf{z}$.
    \item If $2r < d$, there is no such $\mathbf{z}$.
  \end{enumerate}
  How must these statements be modified if $k$ is $2$ or $1$?\kn{Two spheres of radius $r$ about $\mathbf{x}$ and $\mathbf{y}$: they meet
  in a sphere of dimension $k-2$, a point, or nothing. Algebraically, start
  from the midpoint $\mathbf{m} = (\mathbf{x}+\mathbf{y})/2$ and look for
  $\mathbf{z} = \mathbf{m} + \mathbf{v}$ with $\mathbf{v}$ orthogonal to
  $\mathbf{x}-\mathbf{y}$; Exercise 1.17 gives $|\mathbf{v}|^2 = r^2-d^2/4$.
  The condition $k \ge 3$ is what guarantees room for infinitely many such
  $\mathbf{v}$ --- for $k=2$ there are exactly two, for $k=1$ none.}

  \medskip
  \begin{center}
  \begin{tikzpicture}[x=1mm, y=1mm]
    % 2r > d
    \begin{scope}
      \draw[thin] (0,0) circle (9mm); \draw[thin] (10,0) circle (9mm);
      \node[ptfill, minimum size=2pt] at (0,0) {};
      \node[ptfill, minimum size=2pt] at (10,0) {};
      \node[ptfill, minimum size=2.4pt] at (5,7.5) {};
      \node[ptfill, minimum size=2.4pt] at (5,-7.5) {};
      \node[lbl, anchor=north] at (5,-12) {$2r>d$};
    \end{scope}
    % 2r = d
    \begin{scope}[xshift=34mm]
      \draw[thin] (0,0) circle (7mm); \draw[thin] (14,0) circle (7mm);
      \node[ptfill, minimum size=2pt] at (0,0) {};
      \node[ptfill, minimum size=2pt] at (14,0) {};
      \node[ptfill, minimum size=2.4pt] at (7,0) {};
      \node[lbl, anchor=north] at (7,-12) {$2r=d$};
    \end{scope}
    % 2r < d
    \begin{scope}[xshift=72mm]
      \draw[thin] (0,0) circle (5.5mm); \draw[thin] (16,0) circle (5.5mm);
      \node[ptfill, minimum size=2pt] at (0,0) {};
      \node[ptfill, minimum size=2pt] at (16,0) {};
      \node[lbl, anchor=north] at (8,-12) {$2r<d$};
    \end{scope}
  \end{tikzpicture}
  \end{center}
  \figcap{Drawn in the plane, where the intersection is two points, one, or
  none. In $\R^k$ the first case becomes a sphere of dimension $k-2$ --- a
  circle when $k=3$, and infinite whenever $k\ge3$.}

\item Prove that
\[ |\mathbf{x}+\mathbf{y}|^2 + |\mathbf{x}-\mathbf{y}|^2
   = 2|\mathbf{x}|^2 + 2|\mathbf{y}|^2 \]
if $\mathbf{x} \in \R^k$ and $\mathbf{y} \in \R^k$. Interpret this
geometrically, as a statement about parallelograms.\kn{Expand both sides with
the inner product; the cross terms $\pm 2\,\mathbf{x}\cdot\mathbf{y}$ cancel.
Geometrically: the sum of the squares of a parallelogram's diagonals equals the
sum of the squares of its four sides. This identity characterises norms that
come from an inner product, which is why it fails for the norms of Exercise
1.19 in later chapters.}

\medskip
\begin{center}
\begin{tikzpicture}[x=1mm, y=1mm]
  \coordinate (O) at (0,0);
  \coordinate (X) at (30,3);
  \coordinate (Y) at (10,18);
  \coordinate (S) at (40,21);
  \draw[thin] (O) -- (X) -- (S) -- (Y) -- cycle;
  \draw[vec] (O) -- (S);
  \draw[vec] (Y) -- (X);
  \node[ptfill, minimum size=2.4pt] at (O) {};
  \node[mlbl, anchor=north] at (15,1.2) {$\mathbf{x}$};
  \node[mlbl, anchor=east] at (4,10) {$\mathbf{y}$};
  \node[mlbl, anchor=south east] at (26,14) {$\mathbf{x}+\mathbf{y}$};
  \node[mlbl, anchor=north west] at (22,11.5) {$\mathbf{x}-\mathbf{y}$};
  \node[lbl, anchor=west] at (48,14) {sum of squares};
  \node[lbl, anchor=west] at (48,10.5) {of the diagonals};
  \node[lbl, anchor=west] at (48,6) {$=$ sum of squares};
  \node[lbl, anchor=west] at (48,2.5) {of the four sides};
\end{tikzpicture}
\end{center}
\figcap{The identity is exactly this, and it fails for norms that do not come
from an inner product.}

\item If $k \ge 2$ and $\mathbf{x} \in \R^k$, prove that there exists
$\mathbf{y} \in \R^k$ such that $\mathbf{y} \ne \mathbf{0}$ but
$\mathbf{x}\cdot\mathbf{y} = 0$. Is this also true if $k = 1$?\kn{Exhibit one:
if $\mathbf{x} = (x_1,\dots,x_k)$, try $\mathbf{y} = (-x_2,x_1,0,\dots,0)$, and
check separately what to do when $x_1 = x_2 = 0$. For $k=1$ the statement fails
in general but not always: if $x \ne 0$ then $xy = 0$ forces $y = 0$ by 1.16(b),
so no such $y$ exists --- whereas if $x = 0$ then \emph{every} $y \ne 0$ works.
So the honest answer is that it holds precisely when $x = 0$.}

\item Suppose $\mathbf{a} \in \R^k$, $\mathbf{b} \in \R^k$. Find
$\mathbf{c} \in \R^k$ and $r>0$ such that
\[ |\mathbf{x}-\mathbf{a}| = 2|\mathbf{x}-\mathbf{b}| \]
if and only if $|\mathbf{x}-\mathbf{c}| = r$.
\emph{(Solution: $3\mathbf{c} = 4\mathbf{b}-\mathbf{a}$,
$3r = 2|\mathbf{b}-\mathbf{a}|$.)}\kn{Square both sides, expand with the inner
product, and collect: the $|\mathbf{x}|^2$ terms do not cancel --- they combine
to $-3|\mathbf{x}|^2$ --- so dividing by $-3$ leaves the equation of a sphere.
Completing the square identifies the centre. The locus is an Apollonius sphere;
Rudin supplies the answer so you can check your algebra.}

\medskip
\begin{center}
\begin{tikzpicture}[x=1mm, y=1mm]
  \draw[thin] (30,0) circle (13.3mm);
  \node[ptfill, minimum size=2.4pt] at (10,0) {};
  \node[ptfill, minimum size=2.4pt] at (25,0) {};
  \node[ptfill, minimum size=2pt] at (30,0) {};
  \node[mlbl, anchor=north] at (10,-1.6) {$\mathbf{a}$};
  \node[mlbl, anchor=north] at (25,-1.6) {$\mathbf{b}$};
  \node[mlbl, anchor=south] at (30,1.4) {$\mathbf{c}$};
  \coordinate (P) at (39.4,9.4);
  \node[ptfill, minimum size=2.4pt] at (P) {};
  \node[mlbl, anchor=south west] at (P) {$\mathbf{x}$};
  \draw[guide] (10,0) -- (P);
  \draw[guide] (25,0) -- (P);
  \draw[thin, -{Stealth[length=3.5pt]}] (30,0) -- (P);
  \node[lbl, anchor=west] at (48,10) {every $\mathbf{x}$ with};
  \node[lbl, anchor=west] at (48,6.5) {$|\mathbf{x}-\mathbf{a}|=2|\mathbf{x}-\mathbf{b}|$};
  \node[lbl, anchor=west] at (48,2) {lies on one sphere};
\end{tikzpicture}
\end{center}
\figcap{The locus of points whose distances to two fixed points keep a constant
ratio is a sphere --- an Apollonius sphere --- not, as one might guess, a plane.
It is a plane only when the ratio is $1$.}

\item With reference to the Appendix, suppose that property (III) were omitted
from the definition of a cut. Keep the same definitions of order and addition.
Show that the resulting ordered set has the least-upper-bound property, that
addition satisfies axioms (A1) to (A4) (with a slightly different
zero-element!), but that (A5) fails.

\end{enumerate}

\begin{deeper}[On Exercise 1.20]
The exercise that explains a definition rather than applying it. Without (III)
each rational $r$ has \emph{two} representatives, $\{p<r\}$ and $\{p \le r\}$;
the parenthetical hint is that the identity is now $\{p \le 0\}$, not
$\{p<0\}$. Inverses are what break. Take $\alpha = \{p<0\}$, which has no
largest member: then $\alpha+\beta$ has no largest member either, whatever
$\beta$ is, so it can never equal the new zero --- which has largest member
$0$. Compare the note at Step 1.
\end{deeper}

